 \chapter{线性微分方程}\label{ch:b1:diffeq}

微分方程已在高中分册中初见，本章则以
完整证明和更一般的形式讨论：一阶变系数
线性方程（可用常数变易法完全求解）与
二阶常系数线性方程，即振动的模型。两种情形
具有相同的结构，
即 \emph{通解 $=$ 一个特解
$+$ 齐次方程的通解}。

\section{一阶线性方程}

\begin{definition}\label{def:b1:diffeq:linear1}
设 $I$ 为区间，$a, b \colon I \to \R$（或 $\C$）都是
连续的。方程
\[
(E)\colon\quad y' + a(x)\,y = b(x),
\]
在未知的可导函数 $y \colon I \to \R$（或 $\C$）中，
是一个\emph{一阶线性微分
方程}\index{微分方程！一阶线性}。
方程 $(H)\colon y' + a(x) y = 0$ 是它的
\emph{齐次}方程。
\end{definition}

\begin{theorem}[求解齐次方程]\label{thm:b1:diffeq:homogeneous1}
设 $A$ 为 $a$ 在 $I$ 上的一个原函数（存在性见
\cref{ch:b1:integration}）。$(H)$ 在 $I$ 上的解恰为
下列函数：
\[
y(x) = \lambda\, \eu^{-A(x)}, \qquad \lambda \in \R \text{ (或 } \C).
\]
\end{theorem}

\begin{proof}
这些函数都是解：$y' = -\lambda A' \eu^{-A} = -a y$。
反之，设 $y$ 满足 $(H)$，并令 $z(x) = y(x)\, \eu^{A(x)}$。
于是
\[
z' = y' \eu^{A} + y\, a\, \eu^{A} = (y' + ay)\, \eu^{A} = 0 ,
\]
因此 $z$ 在区间 $I$ 上是常数，记 $z = \lambda$，于是 $y =
\lambda \eu^{-A}$。（注意逻辑：
没有解被遗漏，因为 \emph{每个}解都已写成所预告的形式。）
\end{proof}

\begin{example}[变系数的齐次方程]\label{ex:b1:diffeq:varcoeff}
求 $y' + (\cos x)\,y = 0$ 在 $\R$ 上的解。$a(x) =
\cos x$ 的一个原函数是 $A(x) = \sin x$，故解为
\[
y(x) = \lambda\,\eu^{-\sin x}, \qquad \lambda \in \R .
\]
两点说明。每个解都是周期函数（周期为 $2\pi$），且
除 $\lambda = 0$ 外永不为零 --- $\lambda$ 的符号
就是 $y$ 的符号，且始终保持不变，因为指数函数不会穿过零。
满足 $y(0) = y_0$ 的解为 $y_0\eu^{-\sin x}$：
即在每个初始点处，曲线族中恰有一条经过该点，
这正是\cref{thm:b1:diffeq:voc}~(2) 所描述的一维图景。
\end{example}

\begin{theorem}[常数变易法；柯西问题]\label{thm:b1:diffeq:voc}
沿用上述记号：
\begin{enumerate}
  \item $(E)$ 在 $I$ 上的解恰为
        \[
        y(x) = \Bigl(\lambda + \int_{x_0}^{x} b(t)\, \eu^{A(t)} \dd t
        \Bigr)\, \eu^{-A(x)},
        \qquad \lambda \in \R,
        \]
        其中 $x_0 \in I$ 固定。等价地，
        即 $(H)$ 的通解与 $(E)$ 的一个特解之和。
  \item 对任意 $x_0 \in I$ 及 $y_0$，\emph{柯西
        问题}\index{柯西问题}“$(E)$ 且 $y(x_0) = y_0$”
        在 $I$ 上恰有一个解。
\end{enumerate}
\end{theorem}

\begin{proof}
(1) 采用名为\emph{常数变易
法}\index{常数变易法}的方法，寻找形如
$y = \mu(x)\, \eu^{-A(x)}$ 的解（$\mu$ 可导）——这并不损失一般性，
因为 $I$ 上的每个函数都可如此写出，
即 $\mu = y\,\eu^{A}$。代入得
\[
y' + ay = \mu' \eu^{-A} - \mu a \eu^{-A} + a\mu\eu^{-A}
= \mu'\, \eu^{-A},
\]
于是 $y$ 满足 $(E)$ 当且仅当 $\mu'(x) = b(x)\,\eu^{A(x)}$，
亦即 $\mu(x) = \lambda + \int_{x_0}^x b(t)\eu^{A(t)}\dd t$，
其中 $\lambda$ 为某常数（同一区间上连续函数的两个
原函数相差一个常数）。

(2) 在上式中，$y(x_0) = \lambda\,\eu^{-A(x_0)}$，条件
$y(x_0) = y_0$ 唯一确定 $\lambda = y_0 \eu^{A(x_0)}$。
\end{proof}

\begin{example}\label{ex:b1:diffeq:order1}
求 $y' + \dfrac{y}{x} = x^2$ 在 $I = \intoo{0}{+\infty}$ 上的解。这里
$a(x) = \frac 1x$，$A(x) = \ln x$，$\eu^{-A(x)} = \frac 1x$。
齐次方程的解为 $\frac{\lambda}{x}$。用常数变易法：
$\mu'(x) = x^2 \cdot x = x^3$，于是 $\mu = \frac{x^4}{4} + \lambda$，从而
\[
y(x) = \frac{x^3}{4} + \frac{\lambda}{x}, \qquad \lambda \in \R .
\]
由初值条件 $y(1) = 0$ 得 $\lambda = -\frac14$。
\emph{验算：}$y' + \frac yx = \frac{3x^2}{4} - \frac{\lambda}{x^2} +
\frac{x^2}{4} + \frac{\lambda}{x^2} = x^2$。
\end{example}

\begin{example}[猜测胜过积分]\label{ex:b1:diffeq:guess}
求 $y' + 2x\,y = x$ 在 $\R$ 上的解。常数变易法可用（$A =
x^2$，$\mu' = x\,\eu^{x^2}$，$\mu = \frac12\eu^{x^2} + \lambda$），
但注意到\emph{常值}函数 $y_p = \frac12$ 满足方程
（$0 + 2x\cdot\frac12 = x$），这样求解更快。
结合齐次解 $\lambda\,\eu^{-x^2}$，得
\[
y(x) = \frac12 + \lambda\,\eu^{-x^2}, \qquad \lambda \in \R .
\]
每个解都极快地收敛到 $\frac12$，当 $x \to
\pm\infty$ 时，常值特解是一个
\emph{平衡点}，所有解都趋向于它。要点是：
在动用一般方法之前，花十秒钟寻找一个
显然的特解（常值函数、单项式或右端项的倍数）；
随后由结构定理即可完成余下的求解。
\end{example}

\begin{remark}[区间不可忽视]
定理在一个\emph{区间}上成立，此处 $a$ 与 $b$
连续。对于 $y' + \frac yx = 0$，在 $\R^*$ 上，其解为
$\frac{\lambda}{x}$（在 $\intoo{0}{+\infty}$ 上）与
$\frac{\mu}{x}$（在 $\intoo{-\infty}{0}$ 上），其中\emph{独立}的
常数：没有理由让同一个公式
跨过 $0$ 处的奇点而拼接在一起。
\end{remark}

\begin{example}[复右端项，两个实值解]\label{ex:b1:diffeq:complexrhs}
一次求解 $y' - y = \cos x$ 与 $y' - y = \sin x$。在
$\C$ 中考虑右端项 $\eu^{\iu x}$：试取 $y_p =
c\,\eu^{\iu x}$ 得 $c(\iu - 1)\eu^{\iu x} = \eu^{\iu x}$，于是
\[
c = \frac1{\iu - 1} = \frac{-1 - \iu}2,
\qquad
y_p = -\frac{(1 + \iu)(\cos x + \iu\sin x)}2
= \frac{\sin x - \cos x}2 + \iu\,\frac{-\sin x - \cos x}2 .
\]
由于方程的系数为实数，实部与虚部
分离：$\frac{\sin x - \cos x}2$ 是 $y' - y = \cos x$ 的解，
而 $-\frac{\sin x + \cos x}2$ 是 $y' - y = \sin x$ 的解（验证
以前者为例：导数为 $\frac{\cos x + \sin x}2$，减去函数本身
即得 $\cos x$）。一次复数计算
取代了两次常数变易法 --- 这种节省与
\cref{met:b1:diffeq:particular} 对二阶情形作了系统化，
也是\cref{ch:b1:complex}中反复出现的收益。
\end{example}

\section{二阶常系数线性微分方程}

\begin{definition}\label{def:b1:diffeq:linear2}
设 $a, b \in \R$，$f \colon I \to \R$ 连续，方程
\[
(E)\colon\quad y'' + a\,y' + b\,y = f(x)
\]
称为\emph{二阶常系数线性
微分方程}\index{微分方程！二阶线性}；
$(H)\colon y'' + ay' + by = 0$ 为其齐次方程，而
$\chi(r) = r^2 + ar + b$ 为其\emph{特征
多项式}\index{特征多项式}。
\end{definition}

\begin{theorem}[齐次解]\label{thm:b1:diffeq:homogeneous2}
设 $\Delta = a^2 - 4b$ 为 $\chi$ 的判别式。
$(H)$ 在 $\R$ 上的实值解为：
\begin{enumerate}
  \item 若 $\Delta > 0$，记 $r_1 \neq r_2$ 为两个实根：
        $\;y = \lambda\, \eu^{r_1 x} + \mu\, \eu^{r_2 x}$；
  \item 若 $\Delta = 0$，记 $r_0$ 为重根：
        $\;y = (\lambda + \mu x)\, \eu^{r_0 x}$；
  \item 若 $\Delta < 0$，根为 $\alpha \pm \iu\omega$
        （$\omega > 0$）：
        $\;y = \eu^{\alpha x} (\lambda \cos\omega x + \mu
        \sin\omega x)$；
\end{enumerate}
各情形下 $(\lambda, \mu)$ 均取遍 $\R^2$。
\end{theorem}

\begin{proof}
首先注意，对 $r \in \C$，$x \mapsto \eu^{rx}$ 是 $(H)$ 的
解，当且仅当 $\chi(r) = 0$（代入：$(r^2 + ar + b)\eu^{rx} =
0$）。这就是指数函数为何是自然的首选猜测：
求导作用于 $\eu^{rx}$，相当于乘以数
$r$，于是微分方程化为数值方程
$\chi(r) = 0$ --- 整个分析问题归结为
求一个二次多项式的根。

关键一步是作未知函数的变换以降低阶数。设 $r$
为 $\chi$ 的一个（可能为复数的）根，并令 $y = z\, \eu^{rx}$，
这样做不失一般性。于是
\[
y'' + ay' + by
= \bigl(z'' + (2r + a) z' + \chi(r) z\bigr)\eu^{rx}
= \bigl(z'' + (2r + a) z'\bigr)\eu^{rx},
\]
于是 $(H)$ 化为\emph{一阶}方程 $u' + (2r + a) u = 0$
其中 $u = z'$。

\emph{情形 $\Delta \neq 0$：}取 $r = r_1$；则 $2r_1 + a = r_1 -
r_2$（因 $r_1 + r_2 = -a$）。由
\cref{thm:b1:diffeq:homogeneous1}可知，$z' = c\,\eu^{(r_2 - r_1)x}$，
对某常数 $c$；在 $\R$ 上积分，$z = \mu\, \eu^{(r_2 - r_1)x}
+ \lambda$，其中 $\mu = \frac{c}{r_2 - r_1}$，于是 $y = z\,
\eu^{r_1 x} = \lambda\,\eu^{r_1x} + \mu\,\eu^{r_2x}$。
当 $\Delta < 0$ 时，根为 $\alpha \pm \iu\omega$，而
复值解为 $y = c_1\eu^{(\alpha+\iu\omega)x} +
c_2\eu^{(\alpha-\iu\omega)x}$，其中 $c_1, c_2 \in \C$。哪些
为实值函数？由于 $\conj{\eu^{(\alpha+\iu\omega)x}} =
\eu^{(\alpha-\iu\omega)x}$，$y$ 的共轭复数为 $\conj{c_2}\,
\eu^{(\alpha+\iu\omega)x} + \conj{c_1}\,\eu^{(\alpha-\iu\omega)x}$，
而 $y = \conj y$ 对所有 $x$ 成立，推出 $c_2 = \conj{c_1}$（两个
指数函数线性无关：在两点求值，或
在 $x=0$ 处比较（除以 $\eu^{\alpha x}$ 之后）。记 $c_1
= \frac{\lambda - \iu\mu}2$，其中 $\lambda, \mu$ 为实数：
\[
y = 2\,\Re\Bigl(c_1\,\eu^{\alpha x}(\cos\omega x +
\iu\sin\omega x)\Bigr)
= \eu^{\alpha x}\bigl(\lambda\cos\omega x + \mu\sin\omega
x\bigr),
\]
反之，每个这样的函数都是解（实系数方程复值解的
实部）。实解空间
与所述一致。

\emph{情形 $\Delta = 0$：}$r = r_0$，$2r_0 + a = 0$，故 $z'' = 0$：
$z = \lambda + \mu x$，从而 $y = (\lambda + \mu x)\eu^{r_0 x}$。
\end{proof}

\begin{example}[一个柯西问题，从头到尾]\label{ex:b1:diffeq:cauchy2}
求解 $y'' - 3y' + 2y = 0$，且 $y(0) = 0$，$y'(0) = 1$。
特征多项式 $r^2 - 3r + 2 = (r - 1)(r - 2)$ 有
两个实根 $1$ 与 $2$，通解为 $y = \lambda\eu^{x} +
\mu\eu^{2x}$。由两个条件得到线性方程组
\[
\lambda + \mu = 0, \qquad \lambda + 2\mu = 1 ,
\]
于是 $\mu = 1$，$\lambda = -1$：
\[
y(x) = \eu^{2x} - \eu^{x} .
\]
验证：$y(0) = 0$；$y' = 2\eu^{2x} - \eu^x$ 满足 $y'(0) = 1$；且
$y'' - 3y' + 2y = (4 - 6 + 2)\eu^{2x} + (-1 + 3 - 2)\eu^x = 0$。
注意答案的形状：在 $-\infty$ 附近，慢模式
$-\eu^x$ 占主导；在 $+\infty$ 附近则由快模式 $\eu^{2x}$ 主导。
将解看作具有不同衰减或增长速率的
模式的叠加是一个值得养成的习惯 --- 这正是
周末问题中瞬态与稳态分解的组织方式。
\end{example}

\begin{omfigure}
\begin{tikzpicture}[xscale=1.35, yscale=1.55]
  \draw[->, thin] (0,0) -- (6.4,0) node[below] {$x$};
  \draw[->, thin] (0,-0.85) -- (0,1.25) node[left] {$y$};
  \draw[omDef, thick, domain=0:6.2, samples=400, smooth]
    plot (\x, {exp(-0.35*\x)*cos(4*\x r)});
  \draw[omThm, thick, domain=0:6.2, samples=120, smooth]
    plot (\x, {(1+1.8*\x)*exp(-1.1*\x)});
  \draw[omProp, thick, domain=0:6.2, samples=120, smooth]
    plot (\x, {1.15*exp(-0.45*\x) - 0.15*exp(-2.5*\x)});
  \node[omDef, font=\small] at (2.1,-0.62)
    {$\Delta < 0$：振荡};
  \node[omThm, font=\small] at (0.95,1.1) {$\Delta = 0$};
  \node[omProp, font=\small] at (4.9,0.38)
    {$\Delta > 0$：无振荡};
\end{tikzpicture}

{\small 方程 $y'' + ay' + by = 0$ 的三种情形，解均衰减：
阻尼振荡（复根）、临界回归
（二重根）、过阻尼衰减（两个实根）。哪种情形出现，
由 $\Delta = a^2 - 4b$ 的符号直接读出 ---
在求解之前即可判定。}
\end{omfigure}

\begin{theorem}[结构与柯西问题]\label{thm:b1:diffeq:structure2}
\begin{enumerate}
  \item 若 $y_p$ 是 $(E)$ 的一个特解，则
        $(E)$ 的解恰为 $y_p + y_h$，其中 $y_h$ 取遍
        $(H)$ 的解。
  \item （叠加原理）若 $y_1$ 满足 $y'' + ay' + by = f_1$，且
        $y_2$ 满足 $y'' + ay' + by = f_2$，则 $y_1 + y_2$ 满足
        右端为 $f_1 + f_2$ 的方程。
  \item 对所有 $x_0 \in I$ 及 $(y_0, y_0')$，柯西问题
        “$(E)$，$y(x_0) = y_0$，$y'(x_0) = y_0'$”恰有一个
        在 $I$ 上的解。\emph{（存在性由一个特解保证；
        唯一性则完整证明。）}
\end{enumerate}
\end{theorem}

\begin{proof}
(1) $y$ 满足 $(E)$ 当且仅当 $y - y_p$ 满足 $(H)$，这由映射 $y
\mapsto y'' + ay' + by$ 的线性性推出。(2) 同理。

(3) 由 (1)，只需证明常数 $(\lambda,
\mu)$ 总能被唯一地调整，以匹配任意数据 $(y_0, y_0')$。
不妨平移自变量，设 $x_0 = 0$。在
\cref{thm:b1:diffeq:homogeneous2}的情形 (1) 中，$y = \lambda\eu^{r_1x} +
\mu\eu^{r_2x}$ 给出
\[
y(0) = \lambda + \mu, \qquad y'(0) = r_1\lambda + r_2\mu :
\]
一个关于 $(\lambda, \mu)$ 的线性方程组，其行列式为 $r_2 -
r_1 \neq 0$。显式求解得 $\mu = \frac{y_0' -
r_1y_0}{r_2 - r_1}$，$\lambda = y_0 - \mu$，恰有一个
解。在情形 (2) 中，$y(0) = \lambda$，$y'(0) = r_0\lambda +
\mu$，方程组为三角形，行列式为 $1$，解为
$\lambda = y_0$，$\mu = y_0' - r_0y_0$。在情形 (3) 中，$y(0) =
\lambda$，$y'(0) = \alpha\lambda + \omega\mu$，行列式为
$\omega \neq 0$，解为 $\lambda = y_0$，$\mu = \frac{y_0' -
\alpha y_0}\omega$。各情形中，映射 $(\lambda, \mu) \mapsto
(y(x_0), y'(x_0))$ 是线性双射 --- 借助
\cref{ch:b1:linmaps}中的语言，可将此逐例检验压缩为
一句话。
\end{proof}

\begin{method}[$f(x) = P(x)\,\eu^{\gamma x}$ 的特解]\label{met:b1:diffeq:particular}
当右端为 $P(x)\,\eu^{\gamma x}$，其中 $P$ 为
多项式，$\gamma \in \R$（涵盖多项式、指数函数，
经由复数 $\gamma$ 或叠加，亦涵盖 $\cos$ 与 $\sin$）。试求
形如下式的一个特解：
\[
y_p(x) = x^{m}\, Q(x)\, \eu^{\gamma x},
\qquad
m = \text{重数，即 } \gamma \text{ 作为 } \chi
\ (m = 0, 1 \text{ 或 } 2),
\]
其中 $Q$ 为与 $P$ 同次的多项式，其系数
由代入与比较系数确定。对 $f = K\cos\omega x$
（或 $\sin$），取右端为 $K\eu^{\iu\omega x}$ 的方程求解，并取
实部（相应地，虚部）。
\end{method}

\begin{example}[叠加原理的应用]\label{ex:b1:diffeq:superposition}
求 $y'' - y = \eu^{x} + 4$ 在 $\R$ 上的解。齐次方程的特征多项式 $\chi(r) = r^2
- 1$，根为 $\pm1$，故 $y_h = \lambda\eu^x + \mu\eu^{-x}$。将右端拆开，逐项
按方法框处理每一项。
\emph{第一项 $\eu^x$：}此时 $\gamma = 1$ 是
$\chi$ 的单根，故试取 $y_1 = c\,x\,\eu^x$：则 $y_1'' - y_1 = c(x +
2)\eu^x - cx\eu^x = 2c\,\eu^x$，得 $c = \frac12$。
\emph{第二项 $4$：}$\gamma = 0$ 不是根，常数 $y_2 =
-4$ 即可。由叠加原理
（\cref{thm:b1:diffeq:structure2}~(2)）：
\[
y = \frac{x\,\eu^x}2 - 4 + \lambda\,\eu^x + \mu\,\eu^{-x},
\qquad (\lambda, \mu) \in \R^2 .
\]
注意两项要求的形式\emph{不同}（$m = 1$
与 $m = 0$）：重数检验须逐个指数分别施行，这正是
在猜测之前先拆分
右端的全部要义所在。
\end{example}

\begin{example}[重数规则的应用]\label{ex:b1:diffeq:multiplicity}
求 $y'' + y' = x$ 在 $\R$ 上的解。右端为 $P(x)\eu^{0
\cdot x}$，其中 $P(x) = x$，且 $\gamma = 0$ 是 \emph{单}
$\chi(r) = r^2 + r = r(r + 1)$ 的根：故 $m = 1$，正确的
猜测形式为 $y_p = x\,(\alpha x + \beta) = \alpha x^2 + \beta x$，比
$P$ 高一次。代入得：
\[
y_p'' + y_p' = 2\alpha + (2\alpha x + \beta)
= 2\alpha x + (2\alpha + \beta) ,
\]
比较 $x$ 的系数，得 $\alpha = \frac12$，$\beta = -1$：
$y_p = \frac{x^2}2 - x$。通解为 $y = \frac{x^2}2 - x +
\lambda + \mu\,\eu^{-x}$。若忽略重数，猜测 $y_p = \alpha x + \beta$
（即不计重数），代入后得 $y_p'' + y_p'
= \alpha$，为常数 --- 无论 $\alpha, \beta$ 如何选取，都无法与
$x$ 相符，这种失败是结构性的：常数本身已是齐次方程的解，
因而在左端无从显现。
因子 $x^m$ 的存在，正是为了跳出齐次方程的
解空间。
\end{example}

\begin{remark}[第三十二份保险]\label{rem:b1:diffeq:check}
本章每个已解出的方程都以代入检验收尾，
这并非装饰。一个微分方程的计算
往往串联许多小步骤（求原函数、乘积法则、
两个常数），一个符号错误会悄无声息地传播下去；
把最终结果代回方程，
几乎能捕获其中全部错误，代价只是一次求导。
养成这一习惯需分三层：先单独检验\emph{特解}
（齐次部分无论如何都会抵消）；然后在完整解上检验
\emph{初始条件}，若题中含有参数，
则再检验一个\emph{退化值}（一般 $\Omega$ 下的公式
在某取值处能否复现已知结果，
即取 $\Omega = 0$ 时是否成立？）此习惯只需半分钟，却能把
“大概正确”变为“已验证”。
\end{remark}

\begin{remark}[常见误区]\label{rem:b1:diffeq:pitfalls}
\begin{enumerate}
  \item \emph{先化为标准形式。}公式假设方程写成
        $y' + a(x)y = b(x)$，即 $1$ 是 $y'$ 的系数。
        对于 $xy' - 2y = x^3$，两边除以 $x$（在
        不含 $0$ 的区间上）后再确定 $a$ 与 $b$，
        如\cref{exo:b1:diffeq:2}所示。
  \item \emph{每阶一个常数，最后再定。}一阶通解含
        一个常数，二阶通解含两个；初始条件须施加在整个
        \emph{完整解} $y_p + y_h$ 上，而非仅在 $y_h$ 上，
        而在加上 $y_p$ 之前施加条件，
        便是最常见的
        结构性错误。
  \item \emph{注意重数。}特解的试探形式若满足齐次方程，
        则在左端不起作用，
        故 \cref{met:b1:diffeq:particular} 中的因子 $x^m$
        不可省略
        （见\cref{ex:b1:diffeq:multiplicity}）。
  \item \emph{区间是答案的一部分。}解须定义在
        系数连续的区间上；
        若跨越奇点拼接，可能产生虚假的常数
        （见\cref{exo:b1:diffeq:12}），或破坏唯一性。“在
        $\R^*$ 上求解”意味着两个独立的问题。
\end{enumerate}
\end{remark}

\begin{example}[非共振受迫]\label{ex:b1:diffeq:offresonance}
求 $y'' + 4y = \sin x$ 在 $\R$ 上的解。固有频率为 $2$，
外力频率为 $1$；由于 $\iu$ \emph{不}是 $\chi(r) = r^2
+ 4$，故重数 $m = 0$，一个纯正弦即可。
试取 $y_p = \alpha\sin x$（无需余弦项：方程中没有
$y'$ 项，且 $\sin$ 求导后仍为 $\sin$ 型）：
\[
y_p'' + 4y_p = -\alpha\sin x + 4\alpha\sin x = 3\alpha\sin x ,
\]
故 $\alpha = \frac13$，通解为
\[
y = \frac{\sin x}3 + \lambda\cos 2x + \mu\sin 2x .
\]
所有解都有界：它是两个振动的叠加，
频率分别为 $1$（受迫）与 $2$（固有）。对照
下一例中，外力\emph{恰为}固有频率，
会彻底改变解的形态。
\end{example}

\begin{example}[受迫振动]\label{ex:b1:diffeq:oscillation}
求解 $y'' + y = \cos x$，$y(0) = 0$，$y'(0) = 0$。

\emph{齐次部分：}$\chi(r) = r^2 + 1$，根为 $\pm\iu$，$y_h =
\lambda\cos x + \mu \sin x$.

\emph{特解：}右端为 $\Re(\eu^{\iu x})$，其中 $\gamma =
\iu$ 为 $\chi$ 的单根，试取 $z_p = c\, x\, \eu^{\iu x}$
（$c \in \C$）。则 $z_p'' + z_p = c\,(2\iu)\eu^{\iu x}$，其值
等于 $\eu^{\iu x}$，当 $c = \frac{1}{2\iu} = -\frac\iu2$ 时成立。于是 $z_p =
-\frac{\iu}{2} x (\cos x + \iu \sin x)$ 且 $y_p = \Re(z_p) =
\frac{x \sin x}{2}$.

\emph{通解：}$y = \frac{x\sin x}{2} + \lambda\cos x +
\mu\sin x$。条件：$y(0) = \lambda = 0$；$y' = \frac{\sin x +
x\cos x}{2} + \mu\cos x$，故 $y'(0) = \mu = 0$。答案：$y =
\frac{x\sin x}{2}$ --- 一个振幅线性增长的振动：
\emph{共振}\index{共振}现象，由
以固有频率驱动系统所引起。
\end{example}

\begin{omfigure}
\begin{tikzpicture}
\begin{axis}[omaxis, width=11cm, height=5.2cm,
    xmin=0, xmax=25, ymin=-13, ymax=13,
    xtick={0,5,10,15,20,25}, ytick=\empty]
  \addplot[omThm, thick, domain=0:25, samples=400] {x*sin(deg(x))/2};
  \addplot[dashed, gray, domain=0:25] {x/2};
  \addplot[dashed, gray, domain=0:25] {-x/2};
\end{axis}
\end{tikzpicture}

{\small 共振：$y = \frac{x \sin x}{2}$ 是方程 $y'' + y =
\cos x$ 的解，它在直线 $y = \pm\frac x2$（虚线）之间振荡，
振幅不断增大。}
\end{omfigure}

\begin{remark}[插曲：线性即几何]\label{rem:b1:diffeq:interlude}
回顾本章每个解集的形状：
一个特解，加上一族齐次解构成的空间，其中含
一个自由常数（一阶）或两个（二阶）。关于
线性代数各章（\cref{ch:b1:vspaces,ch:b1:findim,ch:b1:linmaps}）
将提供精确的术语：
映射 $L(y) = y'' + ay' + by$ 是\emph{线性}的，其
齐次解构成 $L$ 的\emph{核}，
一个向量空间，其\emph{维数}等于方程的阶 ---
这才是“每阶一个常数”的确切含义 --- 且
方程 $L(y) = f$ 的解集是一个\emph{仿射子空间}，即
核的一个平移。即便是柯西映射 $(\lambda, \mu)
\mapsto (y(x_0), y'(x_0))$ 在\cref{thm:b1:diffeq:structure2}中，
是两平面之间的一个线性双射，亦即一个可逆的 $2
\times 2$ 线性方程组（见\cref{ch:b1:matrices}）。本章
内容无需重做，只需改个名字，
而这一改名正是为线性代数所做的最好热身：
那里的每个抽象定义，在此早已
派上用场。
\end{remark}

\begin{remark}[本章的应用]\label{rem:b1:diffeq:whereused}
结构定理 --- $(E)$ 的解构成“一个特解
加上 $(H)$ 的解” --- 是这一模式的首次出现，
这一模式将在\cref{ch:b1:vspaces,ch:b1:linmaps}中命名：
$(H)$ 的解集是线性映射 $y
\mapsto y'' + ay' + by$ 的\emph{核}，而 $(E)$ 的解集则是
它的一个仿射平移。特征多项式又以
一个矩阵的特征多项式的形式出现在\cref{ch:b1:matrices}中：
二阶方程不过是伪装起来的 $2 \times 2$ 一阶方程组，
这一观点由本科二年级分册系统化。
常数变易法所需的积分由
\cref{ch:b1:integration}提供，下面的周末问题 --- 
受迫阻尼振子 --- 是科学中每一个
振动问题的典型模型，从电路直至
悬索桥。
\end{remark}

\section{练习}

\begin{exercise}[$\star$]\label{exo:b1:diffeq:1}
在 $\R$ 上求解：$\;y' + 2y = \eu^{3x}$；再求柯西问题 $y(0) =
1$。
\end{exercise}

\begin{exercise}[$\star$]\label{exo:b1:diffeq:2}
在 $\intoo{0}{+\infty}$ 上求解：$\;x y' - 2y = x^3$ \emph{（先
将方程化为标准形式）}。
\end{exercise}

\begin{exercise}[$\star$]\label{exo:b1:diffeq:3}
在 $\R$ 上求解，并给出实通解：
$\;y'' - 3y' + 2y = 0$；$\;y'' + 4y' + 4y = 0$；$\;y'' - 2y' + 5y =
0$。
\end{exercise}

\begin{exercise}[$\star$]\label{exo:b1:diffeq:4}
求解 $y'' - y = x^2$，定义域为 $\R$；再求柯西问题 $y(0) = 0$，
$y'(0) = 1$。
\end{exercise}

\begin{exercise}[$\star\star$]\label{exo:b1:diffeq:5}
在 $\intoo{-\frac\pi2}{\frac\pi2}$ 上求解：
$\;y' + y\tan x = \sin 2x$。
\end{exercise}

\begin{exercise}[$\star\star$]\label{exo:b1:diffeq:6}
求解 $y'' - 4y' + 3y = (2x + 1)\,\eu^{x}$，定义域为 $\R$。\emph{（注意
重数：$1$ 是否为特征多项式的根？）}
\end{exercise}

\begin{exercise}[$\star\star$]\label{exo:b1:diffeq:7}
求解 $y'' + 4y = \sin 2x + x$，定义域为 $\R$。\emph{（叠加原理；
对每个右端项分别处理）}。
\end{exercise}

\begin{exercise}[$\star\star$]\label{exo:b1:diffeq:8}
一杯温度为 $T_0 = 80\,^\circ$C 的咖啡放在
温度为 $20\,^\circ$C 的房间中。牛顿冷却定律指出 $T' = -k\,(T - 20)$，其中
$k > 0$。求 $T(t)$，并已知
咖啡降至 $50\,^\circ$C 时已过 $10$ 分钟，求其降至
$25\,^\circ$C 所需的时间。
\end{exercise}

\begin{exercise}[$\star\star\star$]\label{exo:b1:diffeq:9}
（阻尼振子）对 $\varepsilon \geq 0$，考虑
$y'' + 2\varepsilon y' + y = 0$。
\begin{enumerate}
  \item 分别求解 $\varepsilon \in \intco{0}{1}$、$\varepsilon = 1$
        及 $\varepsilon > 1$ 的情形。
  \item 证明：当 $\varepsilon > 0$ 时，每个解都趋于 $0$
        （在 $+\infty$ 处），而当 $\varepsilon = 0$ 时
        非零解则不然。
  \item 对 $\varepsilon \in \intoo{0}{1}$，证明非零解的
        零点等间距分布，间距为
        $\frac{\pi}{\sqrt{1 - \varepsilon^2}}$。
\end{enumerate}
\end{exercise}

\begin{exercise}[$\star\star\star$]\label{exo:b1:diffeq:10}
求所有二阶可导的函数 $f \colon \R \to \R$，使
以下等式成立：
\[
\forall x, y \in \R, \qquad f(x + y) + f(x - y) = 2 f(x) f(y),
\]
且 $f(0) \ne 0$，$f$ 不是常值函数。
\emph{提示：固定 $y$，对 $x$ 两次求导，并在 $0$ 处取值；
证明 $f(0) = 1$，且 $f'' = c f$（$c$ 为某常数）；
再依 $c$ 的符号求解，并验证哪些解满足
函数方程。}
\end{exercise}

\begin{exercise}[$\star\star$]\label{exo:b1:diffeq:11}
（欧拉方程）求解 $x^2 y'' - x y' + y = 0$，定义域为
$\intoo{0}{+\infty}$。\emph{提示：令 $z(t) = y(\eu^t)$，即\
代入 $x = \eu^t$，并证明 $z$ 满足一个线性
常系数方程。}
\end{exercise}

\begin{exercise}[$\star\star\star$]\label{exo:b1:diffeq:12}
考虑整个实直线上的方程 $x\,y' = 2y$，其中
未知的可导函数为 $y \colon \R \to \R$。
\begin{enumerate}
  \item 分别在 $\intoo{0}{+\infty}$ 与 $\intoo{-\infty}{0}$ 上求解。
  \item 证明：对\emph{任意}常数 $a, b \in \R$，
        函数等于 $ax^2$（$x \geq 0$ 时），等于 $bx^2$
        （$x < 0$ 时），它在 $\R$ 上可导，并
        处处满足该方程。
  \item 由此得出 $\R$ 上的解集是一个
        双参数族，并解释这为何不与
        \cref{thm:b1:diffeq:voc}中的唯一性相矛盾。
\end{enumerate}
\end{exercise}

\section{问题：受迫阻尼振子}

\begin{problem}\label{pb:b1:diffeq:1}
一个方程既描述粘性介质中弹簧上的质量块，
也描述 RLC 电路中的电荷，以及风中摇摆的建筑物：
\[
(E_\Omega)\colon\quad
x'' + 2\lambda x' + \omega_0^2\,x = A\cos(\Omega t),
\]
其中 $\lambda \geq 0$ 为阻尼系数，$\omega_0 > 0$ 为固有
频率，$A > 0$、$\Omega > 0$ 分别为
强迫项的振幅与频率。本题将完整刻画其行为：瞬态的衰减，
唯一的周期稳态、共振曲线
及其尖锐程度（\emph{品质因数}），无阻尼情形的拍，
以及维持振动的
能量平衡。\index{共振}\index{振子}
除非另有说明，设 $0 < \lambda < \omega_0$（欠阻尼
情形），并记 $\omega_d = \sqrt{\omega_0^2 - \lambda^2}$。

\medskip
\noindent\textbf{第一部分 --- 自由振子。}此时 $A = 0$。
\begin{enumerate}
  \item 求齐次方程 $(H)$ 在 $0 < \lambda <
        \omega_0$ 及 $\lambda = 0$ 时的解。（$\lambda
        \geq \omega_0$ 的情形已在\cref{exo:b1:diffeq:9}中讨论，
        请引用其结论。）
  \item 证明：对每个 $\lambda > 0$，$(H)$ 的一切解都
        以 $0$ 为其在 $+\infty$ 处的极限 --- 三种情形皆然。
  \item 定义能量 $\mathcal E(t) = \frac12 x'(t)^2 +
        \frac12\omega_0^2\,x(t)^2$（沿 $(H)$ 的解）。证明
        $\mathcal E'(t) = -2\lambda\,x'(t)^2 \leq 0$，并由此（无需求解）推出，
        柯西问题
        “$(H)$，$x(t_0) = x'(t_0) = 0$”只有
        零解，对每个 $\lambda \geq 0$ 均成立。
  \item 当 $0 < \lambda < \omega_0$ 时，将非零解写成
        $x(t) = R\,\eu^{-\lambda t}\cos(\omega_d t - \varphi)$，并记
        伪周期 $T_d = \frac{2\pi}{\omega_d}$。
        证明 $x(t + T_d) = \eu^{-\lambda T_d}\,x(t)$，即每次
        摆动都是前一次摆动缩小为常数因子
        $\eu^{-\delta}$ 倍所得，其中 $\delta = \frac{2\pi\lambda}{\omega_d}$
        （称为\emph{对数衰减率}）。求 $\delta$，
        取 $\omega_0 = 1$、$\lambda = 0.1$。
  \item 定义\emph{品质因数} $Q =
        \dfrac{\omega_0}{2\lambda}$\index{品质因数}。证明：
        经过时间 $\frac1\lambda$（振幅的一个 $\eu$ 倍衰减）后，
        振子已完成
        $\frac{\omega_d}{2\pi\lambda}$ 个伪周期，当阻尼较弱
        （$\lambda \ll \omega_0$）时，这一数值近似为
        $\frac Q\pi$。在相差 $\pi$ 的意义下，品质因数
        度量了振幅衰减 $\eu$ 倍之前
        所经历的振动次数。
\end{enumerate}

\medskip
\noindent\textbf{第二部分 --- 稳态。}此时 $A > 0$ 且
$\lambda > 0$。
\begin{enumerate}[resume]
  \item 寻找特解，形式为
        $z\,\eu^{\iu\Omega t}$ 的实部，其中 $z \in \C$
        （见\cref{met:b1:diffeq:particular}）。证明可取
        使下式成立：
        \[
        z = \frac{A}{\omega_0^2 - \Omega^2 + 2\iu\lambda\Omega} .
        \]
  \item 由此得出振幅--相位形式的稳态解：
        $x_p(t) = R(\Omega)\cos\bigl(\Omega t -
        \varphi(\Omega)\bigr)$，其中
        \[
        R(\Omega) = \frac{A}{\sqrt{(\omega_0^2 - \Omega^2)^2 +
        4\lambda^2\Omega^2}},
        \qquad
        \tan\varphi = \frac{2\lambda\Omega}{\omega_0^2 -
        \Omega^2},
        \quad \varphi \in \intoo0\pi .
        \]
  \item 解释两种极端情形：求
        $R$ 与 $\varphi$ 当 $\Omega \to 0^+$ 时的极限（准静态响应为
        $A/\omega_0^2$，相位为 $0$），以及当 $\Omega \to
        +\infty$ 时的极限（$R \sim A/\Omega^2 \to 0$，相位 $\to \pi$：
        质量块的运动与过快的强迫反相）。
  \item 证明：\emph{每个}$(E_\Omega)$ 的解都等于 $x_p$
        与 $(H)$ 的某个解之和，因此
        收敛于稳态 $x_p$（当 $t \to +\infty$ 时），
        与初始条件无关。瞬态衰减之后，
        振子已不再记得它是如何开始的。
  \item 证明：$x_p$ 是
        $(E_\Omega)$ 的\emph{唯一}周期解。
  \item 完整求解一个柯西问题：对于 $x'' + 2x' + 2x =
        \cos t$，满足 $x(0) = x'(0) = 0$，证明其解为
        \[
        x(t) = \frac{\cos t + 2\sin t}5
        - \eu^{-t}\,\frac{\cos t + 3\sin t}5 ,
        \]
        并分辨瞬态部分与稳态部分。
\end{enumerate}

\medskip
\noindent\textbf{第三部分 --- 共振曲线。}考察
$\Omega \mapsto R(\Omega)$ 在 $\intoo0{+\infty}$ 上的性态。
\begin{enumerate}[resume]
  \item 令 $u = \Omega^2$，$g(u) = (\omega_0^2 - u)^2 +
        4\lambda^2 u$。证明：若 $2\lambda^2 < \omega_0^2$，则 $R$ 在
        \emph{共振频率}处取得严格极大值，
        $\Omega_r = \sqrt{\omega_0^2 - 2\lambda^2}$，其值为
        \[
        R_{\max} = R(\Omega_r)
        = \frac{A}{2\lambda\sqrt{\omega_0^2 - \lambda^2}} .
        \]
  \item 证明 $\dfrac{R_{\max}}{R(0)} = Q\,\bigl(1 -
        \tfrac{\lambda^2}{\omega_0^2}\bigr)^{-1/2} \geq Q$：
        共振时，强迫项被放大的倍数
        （近似地）等于品质因数。
  \item 证明：\emph{速度}振幅 $V(\Omega) =
        \Omega\,R(\Omega)$ 恰在 $\Omega =
        \omega_0$（而非 $\Omega_r$）处取极大值，且该处的相位为
        $\varphi(\omega_0) = \frac\pi2$。在 $\Omega = \omega_0$ 时，
        速度与强迫项恰好同相。
  \item （带宽）精确求解 $g(u) = 2\,g(u_r)$，其中 $u_r =
        \omega_0^2 - 2\lambda^2$，并由此推出：两个频率 $\Omega_\pm$ 满足
        $R = R_{\max}/\sqrt2$，且 $\Omega_+^2 - \Omega_-^2 =
        4\lambda\sqrt{\omega_0^2 - \lambda^2}$；由此可得，
        弱阻尼时带宽为 $\Omega_+ - \Omega_-
        \approx 2\lambda$，即 $Q \approx
        \frac{\omega_0}{\Omega_+ - \Omega_-}$：尖锐的
        共振峰
        对应高 $Q$ 的系统。
  \item 数值刻画：取 $\omega_0 = 1$，$\lambda = 0.05$
        （$Q = 10$），$A = 1$，求 $\Omega_r$、$R_{\max}$、
        静态响应 $R(0)$，以及近似带宽。
  \item 证明：若 $2\lambda^2 \geq \omega_0^2$，则 $R$ 在
        $\intoo0{+\infty}$ 上严格递减：强阻尼系统
        根本不存在共振峰。
\end{enumerate}

\medskip
\noindent\textbf{第四部分 --- 无阻尼：拍与共振。}
此时 $\lambda = 0$。
\begin{enumerate}[resume]
  \item 当 $\Omega \neq \omega_0$ 时，求方程
        $x'' + \omega_0^2 x = A\cos(\Omega t)$ 的通解。
  \item 求柯西问题 $x(0) = x'(0) = 0$ 的解，并将答案
        化为乘积形式：
        \[
        x(t) = \frac{2A}{\omega_0^2 - \Omega^2}\,
        \sin\Bigl(\frac{(\omega_0 - \Omega)t}2\Bigr)
        \sin\Bigl(\frac{(\omega_0 + \Omega)t}2\Bigr) .
        \]
  \item 当 $\Omega$ 接近 $\omega_0$ 时，把乘积看作
        频率为 $\frac{\omega_0 + \Omega}2$ 的快振荡，
        被一个慢包络调制，其频率为
        $\frac{\abs{\omega_0 - \Omega}}2$，这就是\emph{拍}
        \index{拍}。给出包络的周期和
        最大振幅，并说明当 $\Omega
        \to \omega_0$ 时两者如何无限增大。
  \item 固定 $t$，令 $\Omega \to \omega_0$，在第 19 题的
        公式中取极限，证明极限为
        \[
        x_\infty(t) = \frac{A\,t\,\sin(\omega_0 t)}{2\omega_0} ,
        \]
        并直接验证 $x_\infty$ 满足共振方程
        $x'' + \omega_0^2 x = A\cos(\omega_0 t)$，且
        $x(0) = x'(0) = 0$（比较
        \cref{ex:b1:diffeq:oscillation}）：共振是
        越来越慢、越来越大的拍的极限。
  \item 对比共振的两种结局：无阻尼时呈线性增长
        $\frac{At}{2\omega_0}$，而弱阻尼时
        则饱和于 $R_{\max} \approx Q\,\frac{A}{\omega_0^2}$。
        用一句话说明：什么物理机制把前一种
        结局转化为后一种？
\end{enumerate}

\medskip
\noindent\textbf{第五部分 --- 能量平衡与综合。}
\begin{enumerate}[resume]
  \item 在第二部分的稳态中，求
        一个周期 $\frac{2\pi}\Omega$ 上的平均值：(a) 外力输入的功率
        即 $P_{\mathrm{in}}(t) = A\cos(\Omega t)
        \cdot x_p'(t)$，(b) 阻尼耗散的功率
        即 $P_{\mathrm{diss}}(t) = 2\lambda\,x_p'(t)^2$。
        证明两个平均值都等于 $\lambda\,R^2\Omega^2$：外力
        输入的功率恰好被阻尼耗尽 --- 这正是稳态
        能够保持稳定的原因。
  \item 本题究竟用到了：(i) 结构
        \cref{thm:b1:diffeq:structure2}；(ii) 复指数
        方法；(iii) 一项实变量函数研究，
        风格类似 \cref{ch:b1:functions}？
        每项用一句话回答。
  \item 综合：描述 $(E_\Omega)$ 的完整行为图景
        --- 自由与受迫、有阻尼与无阻尼，以及
        $Q$ 作为唯一的无量纲旋钮，调节峰高、
        带宽与暂态寿命 --- 并指出
        内容的延续方向：一阶 $2 \times 2$
        系统（见\cref{ch:b1:matrices}及本科二年级分册）以及
        一般周期外力分解为
        正弦函数之和（傅里叶级数，见本科三年级分册），
        而本题的正弦情形正是其中的
        基本构件。
\end{enumerate}
\end{problem}
